%% % Plantilla de Documento Final según Normas
% 1) ICONTEC NTC 1486
% 2) ICONTEC NTC 5613
% 3) Estilo de citas bibliográficas de IEEE
%
%% % Autor:
% Germán A Holguín L @ Universidad Tecnológica de Pereira
%
%% % Licencia:
% LaTeX Project Public License (LPPL) (C) SEPTIEMBRE DE 2014
% Lo que significa: Puede usarse libremente con fines no comerciales.
%                   Puede modificarse libremente con fines no comerciales.
%                   Puede distribuirse libremente con fines no comerciales.
%                   Se entrega COMO ESTA. No hay garantía ni soporte.
%
%                   No borrar este mensaje y citar la plantilla de forma adecuada.
%
% Versión plantilla : 13 de Septiembre de 2016

%% Copyright 2016 Germán A. Holguín
%
% This work may be distributed and/or modified under the
% conditions of the LaTeX Project Public License, either version 1.3
% of this license or (at your option) any later version.
% The latest version of this license is in
%   http://www.latex-project.org/lppl.txt
% and version 1.3 or later is part of all distributions of LaTeX
% version 2005/12/01 or later.
%
% This work has the LPPL maintenance status `unmaintained'. Provided AS IS.
% 
% The Current Maintainer of this work is Germán A. Holguín.
%
% This work consists of the files DocuFinal.tex, Ejemplo_Referencias.bib, and
% tesisUTP.bst.

% LPPL means : You are free to use it and modify it as you please. Just keep this message and
%              do not try to make money out of it.


\documentclass[12pt,twoside]{article}
\renewcommand{\rmdefault}{cmr}
\renewcommand{\sfdefault}{cmss}
\renewcommand{\ttdefault}{cmtt}
\usepackage[T1]{fontenc}
\usepackage[utf8]{inputenc}
\usepackage[english,spanish]{babel}
\usepackage[letterpaper]{geometry}
\geometry{verbose,nomarginpar,tmargin=3cm,bmargin=3cm,lmargin=3cm,rmargin=3cm}
\setlength{\parskip}{\medskipamount}
\setlength{\parindent}{0pt}
\usepackage{array}
\usepackage{float}
\usepackage{setspace}
\doublespacing
\usepackage[unicode=true,
 bookmarks=true,bookmarksnumbered=false,bookmarksopen=false,
 breaklinks=false,pdfborder={0 0 0},backref=section,colorlinks=true, linkcolor=blue]
 {hyperref}
\hypersetup{pdftitle={Titulo del proyecto sin tildes},
 pdfauthor={Nombres del los autores sin tildes},
 pdfsubject={tema de la investigacion},
 pdfkeywords={palabras claves}}

\makeatletter

%% Los convertidores HTML no saben que es tabularnewline
\providecommand{\tabularnewline}{\\}

% No generar la salida para \date
\date{}
% Utilizar el punto decimal, en vez de la acostumbrada coma del español.
% Esto está de acuerdo con la Norma NTC/ISO 80000-1
\decimalpoint

% Etiqueta de figura y tabla a la izquierda seguida de punto.
\usepackage[justification=justified, singlelinecheck=false, labelsep = period]{caption}

% Centrar títulos de secciones
\usepackage{sectsty}
\sectionfont{\centering}

% Empezar sección en pagina impar
\let\stdsection\section
\renewcommand\section{\cleardoublepage\stdsection}

% Renombrar las secciones de acuerdo con NTC1486
\addto{\captionsspanish}{\renewcommand*{\contentsname}{\centering CONTENIDO}}
\addto{\captionsenglish}{\renewcommand*{\contentsname}{\centering CONTENIDO}}
\addto{\captionsspanish}{\renewcommand*{\listfigurename}{\centering LISTA DE FIGURAS}}
\addto{\captionsenglish}{\renewcommand*{\listfigurename}{\centering LISTA DE FIGURAS}}
\addto{\captionsspanish}{\renewcommand*{\listtablename}{\centering LISTA DE TABLAS}}
\addto{\captionsenglish}{\renewcommand*{\listtablename}{\centering LISTA DE TABLAS}}
\addto{\captionsspanish}{\renewcommand*{\refname}{BIBLIOGRAF\'IA}}
\addto{\captionsenglish}{\renewcommand*{\refname}{BIBLIOGRAF\'IA}}
\addto{\captionsspanish}{ \renewcommand*{\tablename}{Tabla}}
\addto{\captionsenglish}{ \renewcommand*{\tablename}{Tabla}}
\addto{\captionsenglish}{ \renewcommand*{\figurename}{Figura}}

%Adicionar la palabra Pág. encima de los números
\addtocontents{toc}{~\hfill\textbf{pág.}\par}

% Crear un nuevo tipo de sección que sea centrado y sin numeración.
\makeatletter
\newcommand\@csection{\@startsection {section}{1}{\z@}%
  {-3.5ex \@plus -1ex \@minus -.2ex}%
  {2.3ex \@plus.2ex}%
  {\normalfont\Large\bfseries\centering}}
\newcommand{\csection}[1]{%
  \@csection*{#1}%
}
\makeatother

% Centrar la palabra BIBLIOGRAFíA, sin modificar la entrada en la TOC.
\makeatletter
    \renewenvironment{thebibliography}[1]{%
    \csection{\refname} \bigskip
    %\@mkboth{\MakeUppercase\refname}{\MakeUppercase\refname}%
    \list{\@biblabel{\@arabic\c@enumiv}}
            {\settowidth\labelwidth{\@biblabel{#1}}%
    \leftmargin\labelwidth
    \advance\leftmargin\labelsep
           \@openbib@code
    \usecounter{enumiv}%
    \let\p@enumiv\@empty
    \renewcommand\theenumiv{\@arabic\c@enumiv}}%
    \sloppy
    \clubpenalty4000\@clubpenalty \clubpenalty
    \widowpenalty4000%
    \sfcode`\.\@m}
           {\def\@noitemerr{\@latex@warning{Empty `thebibliography' environment}}%
    \endlist}
\makeatother

\makeatother

% Paquete para generar texto ficticio Lorem Ipsum
% se puede borrar cuando ya todo el texto sea reemplazado.
\usepackage{lipsum}

%%%% ------ EL DOCUMENTO EMPIEZA AQUÍ --------------------------------------------
\begin{document}
\noindent \begin{center}
{\footnotesize{}\pdfbookmark[1]{TíTULO}{TITULO}\vspace{1cm}
}\textbf{TÍTULO DEL PROYECTO DE GRADO}
\par\end{center}

\noindent \begin{center}
\vspace{0.7cm}

\par\end{center}

\begin{center}
Nombre Completo Autor 1 \\ Nombre Completo Autor 2
\par\end{center}

\noindent \begin{center}
\vspace{0.7cm}

\par\end{center}

\noindent \begin{center}
Proyecto de grado presentado como requisito parcial\\para aspirar
al título de Ingeniero Electricista
\par\end{center}

\noindent \begin{center}
\vspace{0.7cm}

\par\end{center}

\noindent \begin{center}
Director\\
Nombre del Director del Proyecto\\
\par\end{center}

\noindent \begin{center}
Grupo de Investigación en Gestión de \\
Sistemas Eléctricos, Electrónicos y Automáticos.
\par\end{center}

\noindent \begin{center}
\vspace{0.7cm}

\par\end{center}

\noindent \begin{center}
\textbf{UNIVERSIDAD TECNOLÓGICA DE PEREIRA\\PROGRAMA DE INGENIERÍA
ELÉCTRICA\\PEREIRA\\20XX}
\par\end{center}

\noindent \begin{center}
\pagenumbering{gobble}\cleardoublepage{}
\par\end{center}

\noindent \begin{flushright}
\begin{tabular}{c>{\centering}p{0.45\textwidth}}
 & Nota de Aceptación\tabularnewline
 & \tabularnewline
\cline{2-2} 
 & \tabularnewline
\cline{2-2} 
 & \tabularnewline
\cline{2-2} 
 & \tabularnewline
\cline{2-2} 
 & \tabularnewline
 & \tabularnewline
\cline{2-2} 
 & Firma del Presidente del jurado\tabularnewline
 & \tabularnewline
 & \tabularnewline
 & \tabularnewline
\cline{2-2} 
 & Firma del jurado 1 - Evaluador\tabularnewline
 & \tabularnewline
 & \tabularnewline
 & \tabularnewline
\cline{2-2} 
 & Firma del jurado 2 - Evaluador\tabularnewline
 & \tabularnewline
 & \tabularnewline
 & \tabularnewline
\cline{2-2} 
 & Firma del jurado 3 - Director\tabularnewline
 & \tabularnewline
 & \tabularnewline
\end{tabular}
\par\end{flushright}

Pereira, XX de XXX de 20XX

\cleardoublepage{}

Dedicado a ....

\cleardoublepage{}Agradecimientos .... 

\cleardoublepage{}

\pdfbookmark[1]{CONTENIDO}{CONTENIDO}\tableofcontents{}\listoftables
\listoffigures



\section{\label{sec:introduccion}INTRODUCCIÓN}

\pagenumbering{arabic}
\setcounter{page}{11}

% Las lineas lipsum crean el texto ficticio.
% simplemente reemplace estas lineas por su contenido real.
\lipsum[1-5]

%% Ejemplo de Tabla.
% La diferencia con las figuras, es que llevan el caption debajo.
\begin{table}[H]
\noindent \begin{raggedright}
\caption{Aquí va el título de la Tabla 1.}

\par\end{raggedright}

\noindent \centering{}%
\begin{tabular}{|c|c|c|c|}
\hline 
 & Col 1 & Col 2 & Col 3\tabularnewline
\hline 
\hline 
Row 1 & A &  & \tabularnewline
\hline 
Row 2 &  & B & \tabularnewline
\hline 
Row 3 &  &  & C\tabularnewline
\hline 
\end{tabular}
\end{table}

% Las lineas lipsum crean el texto ficticio.
% simplemente reemplace estas lineas por su contenido real.
\lipsum[1-2]


\subsection{\label{sec:definicion_del_problema}DEFINICIÓN DEL PROBLEMA}
% Las lineas lipsum crean el texto ficticio.
% simplemente reemplace estas lineas por su contenido real.
\lipsum[1-5]


\subsection{\label{sec:justificacion}JUSTIFICACIÓN}
% Las lineas lipsum crean el texto ficticio.
% simplemente reemplace estas lineas por su contenido real.
\lipsum[1-5]


\subsection{\label{sec:objetivos}OBJETIVOS}
% Las lineas lipsum crean el texto ficticio.
% simplemente reemplace estas lineas por su contenido real.
\lipsum[1-1]

\subsubsection{\label{sub:objetivo_general}Objetivo General}
% Las lineas lipsum crean el texto ficticio.
% simplemente reemplace estas lineas por su contenido real.
\lipsum[1-1]


\subsubsection{\label{sub:objetivos_especificos}Objetivos Específicos}
\begin{itemize}
\item Lorem ipsum dolor sit amet, consectetur adipiscing elit.
\item Lorem ipsum dolor sit amet, consectetur adipiscing elit.
\item Lorem ipsum dolor sit amet, consectetur adipiscing elit.
\item Lorem ipsum dolor sit amet, consectetur adipiscing elit.
\end{itemize}

\section{\label{sec:estado_arte}ESTADO DEL ARTE}
% Las lineas lipsum crean el texto ficticio.
% simplemente reemplace estas lineas por su contenido real.
\lipsum[1-5]


\section{\label{sec:marcoteorico}MARCO TEÓRICO}
% Las lineas lipsum crean el texto ficticio.
% simplemente reemplace estas lineas por su contenido real.
\lipsum[1-5]


\section{\label{sec:metodologia}METODOLOGÍA}
% Las lineas lipsum crean el texto ficticio.
% simplemente reemplace estas lineas por su contenido real.
\lipsum[1-5]


\section{\label{sec:experimentos}EXPERIMENTOS Y RESULTADOS}
% Las lineas lipsum crean el texto ficticio.
% simplemente reemplace estas lineas por su contenido real.
\lipsum[1-5]


\section{\label{sec:conclusiones}CONCLUSIONES Y RECOMENDACIONES}
% Las lineas lipsum crean el texto ficticio.
% simplemente reemplace estas lineas por su contenido real.
\lipsum[1-1]

\subsection{CONCLUSIONES}
% Las lineas lipsum crean el texto ficticio.
% simplemente reemplace estas lineas por su contenido real.
\lipsum[1-5]


\subsection{RECOMENDACIONES}
% Las lineas lipsum crean el texto ficticio.
% simplemente reemplace estas lineas por su contenido real.
\lipsum[1-5]

% Ejemplo de cita:
% Las lineas lipsum crean el texto ficticio.
% simplemente reemplace estas lineas por su contenido real.
\lipsum[1-5]~\cite{Ramirezetal2013GuiaIcontec}

\cleardoublepage{}

\bibliographystyle{tesisUTP}
\phantomsection\addcontentsline{toc}{section}{\refname}
% Borrar linea \nocite{*} en el documento real.
% \nocite{*} imprime todas las referencias del archivo .bib, sin importar sin han sido referenciadas o no.  Por eso se debe borrar.
\nocite{*}
\bibliography{Ejemplo_Referencias}

hello

\end{document}

